\originalchapter{8.223作业三}
题源: Exercises 3, 原题23--31, 题面PDF第1--2页; 本章除原题29明确标出的独立修正外, 依据官方Solutions 3中文改写并复算. 第30--31题补充了原式的初值退化与复特征值.

\begin{exercise}
原题23. 对一维谐振子直接求$T(t),U(t)$及一个周期平均, 验证位力定理.
\end{exercise}
\begin{solution}
官方解答中文改写. 取$x=A\cos(\omega t+\phi)$, $\omega^2=k/m$. 动能$T=kA^2\sin^2(\omega t+\phi)/2$, 势能$U=kA^2\cos^2(\omega t+\phi)/2$. 两平方三角函数一个周期的平均都是$1/2$, 所以$\langle T\rangle=\langle U\rangle=kA^2/4$. 二次势的位力定理$2\langle T\rangle=2\langle U\rangle$得到验证.
\end{solution}

\begin{exercise}
原题24. 固定不可穿透球半径$R$, 球外$U=0$, 球内$U=\infty$. 求散射角绝对值大于$\pi/2$的背散射截面.
\end{exercise}
\begin{solution}
官方解答中文改写. 光滑硬球的弹性反射满足入射角等于反射角. 令瞄准距离为$b=R\sin\alpha$, 则偏折角$\Theta=\pi-2\alpha$. 背散射要求$\alpha<\pi/4$, 即$b<R/\sqrt2$, 所以$\sigma_{\rm back}=\pi R^2/2$. 零测度边界的严格不等号不影响截面.
\end{solution}

\begin{exercise}
原题25. 对共振强迫$\ddot x+\omega_0^2x=(F/m)\cos(\omega_0t+\theta)$证明解可写成$a_1\cos(\omega_0t+\phi)+a_2t\sin(\omega_0t+\theta)$, 并用$x(0)=0$, $\dot x(0)=v_0$求常量.
\end{exercise}
\begin{solution}
官方解答中文改写. 特解求导给$\ddot x_p+\omega_0^2x_p=2a_2\omega_0\cos(\omega_0t+\theta)$, 所以$a_2=F/(2m\omega_0)$. 不必通过有歧义的幅相直接求初值, 写齐次解$A\cos\omega_0t+B\sin\omega_0t$, 得$A=0$, $B=(v_0-a_2\sin\theta)/\omega_0$. 因而
\[
x(t)=\frac{v_0-a_2\sin\theta}{\omega_0}\sin\omega_0t+a_2t\sin(\omega_0t+\theta).
\]
若允许$a_1$为有符号系数, 可取$\phi=\pi/2$, $a_1=(a_2\sin\theta-v_0)/\omega_0$. 若要求$a_1\ge0$, 则相位按该系数正负选$\pm\pi/2$; 系数为零时相位无关.
\end{solution}

\begin{exercise}
原题26（双倍权重）. 对原题16的调速器, 求稳定平衡附近的小振动角频率.
\end{exercise}
\begin{solution}
官方解答中文改写, 补充轴线平衡. 有效势为$V_{\rm eff}=-m_1a^2\Omega^2\sin^2\theta-2(m_1+m_2)ga\cos\theta$, 有效质量$2a^2(m_1+2m_2\sin^2\theta)$. 非轴线平衡满足$\cos\theta_0=(m_1+m_2)g/(m_1a\Omega^2)$. 当$\Omega^2>(m_1+m_2)g/(m_1a)$且$0<\theta_0<\pi/2$时存在, 线性化得
\[
\omega^2=
\frac{(m_1+m_2)g\cos\theta_0-m_1a\Omega^2\cos2\theta_0}
{a(m_1+2m_2\sin^2\theta_0)}
=\frac{m_1\Omega^2\sin^2\theta_0}{m_1+2m_2\sin^2\theta_0}>0.
\]
轴线$\theta_0=0$也满足方程, 其$\omega^2=(m_1+m_2)g/(m_1a)-\Omega^2$, 在低转速下稳定; 临界转速频率为零, 需要高阶分析. 轴线姿态中两球重合是理想点质量模型, 真实装置的几何碰撞需另加约束.
\end{solution}

\begin{exercise}
原题27. (a, 双倍权重) 对$m\ddot x=-\lambda\dot x-kx$用指数试探解, 求欠阻尼、临界阻尼、过阻尼三类通解及其指数、频率. (b) 加强迫$C\sin\omega t$, 将解分为暂态和稳态, 写稳态幅相, 求$\omega_0,D,\Gamma,\phi$使$x_{II}=D\sin(\omega t+\phi)/\sqrt{(\omega^2-\omega_0^2)^2+\Gamma^2}$.
\end{exercise}
\begin{solution}
官方解答中文改写. 指数$x=e^{st}$给$ms^2+\lambda s+k=0$. 设$\Lambda=\lambda/(2m)$, $\omega_0^2=k/m$. 当$\lambda^2<4mk$时$x=e^{-\Lambda t}(A\sin\omega_dt+B\cos\omega_dt)$, $\omega_d=\sqrt{\omega_0^2-\Lambda^2}$. 等号时$x=(A+Bt)e^{-\Lambda t}$; 大于时$x=Ae^{s_+t}+Be^{s_-t}$, $s_\pm=(-\lambda\pm\sqrt{\lambda^2-4mk})/(2m)$. 两根之所以合并时要加$t$因子, 可由对指数根求导或将重复根方程逐次积分验证.

(b) 取稳态为$\Im(Xe^{i\omega t})$, $X=(C/m)/(\omega_0^2-\omega^2+i\lambda\omega/m)$. 因而$D=C/m$, $\Gamma=\lambda\omega/m$, $\phi=-\operatorname{atan2}(\lambda\omega/m,\omega_0^2-\omega^2)$. 若$C<0$, 可保留有符号$D$或把$\pi$吸收到相位. 当阻尼正时暂态最终衰减, 相位使用双参数反正切以避免高于共振时落入错误象限.
\end{solution}

\begin{exercise}
原题28. 对$\ddot x+\omega_0^2x=(F/m)\cos(\gamma t+\alpha)$, 非共振解可写成
$x=C\cos(\omega_0t+\kappa)+(F/m)[\cos(\gamma t+\alpha)-\cos(\omega_0t+\alpha)]/(\omega_0^2-\gamma^2)$.
(a) 令$\omega_0=\gamma+\epsilon$, $|\epsilon|\ll\omega_0$, 将分母保留至$\epsilon$一阶, 写题目要求的近共振形式. (b) 在$\epsilon\to0$时得到共振解; 原题随后给Landau记号对应, 本册使用统一记号.
\end{exercise}
\begin{solution}
官方解答中文改写. 因$\omega_0^2-\gamma^2=2\omega_0\epsilon-\epsilon^2$, 一阶幅值近似为
\[
x\simeq C\cos(\omega_0t+\kappa)+\frac{F}{2m\omega_0\epsilon}
[\cos(\omega_0t+\alpha-\epsilon t)-\cos(\omega_0t+\alpha)].
\]
这里保留慢相位$\epsilon t$有助描述拍频, 不将长时间的相位误差忽略. 若取$\epsilon\to0$且$t$固定, 方括号除以$\epsilon$趋于$t\sin(\omega_0t+\alpha)$, 得共振特解$Ft\sin(\omega_0t+\alpha)/(2m\omega_0)$. 非共振公式中的发散分母由分子差抵消, 不能将两个余弦项拆开后各自取极限.
\end{solution}

\begin{exercise}
原题29（双倍权重）. 摆支点水平高频振动$x_s=a\cos\gamma t$, $\gamma\gg\sqrt{g/\ell}$. 求慢运动的稳定平衡位置（水平Kapitza摆）.
\end{exercise}
\begin{solution}
AI独立修正解答. 官方Solutions 3在本题的平均项符号、有效势和稳定性陈述间不一致, 不能把其最终“倒立稳定”结论直接用于水平驱动. 以下按支点位置重新推导, 并额外声明平均近似要求$a/\ell\ll1$.

令$\theta$从竖直向下量起. 位置$x=x_s+\ell\sin\theta$, $y=-\ell\cos\theta$, 完整方程为$\ddot\theta+(g/\ell)\sin\theta=-(\ddot x_s/\ell)\cos\theta=(a\gamma^2/\ell)\cos\gamma t\cos\theta$. 写$\theta=\Theta+\xi$, $\Theta$慢、$\xi$快, 首阶快解$\xi=-(a/\ell)\cos\gamma t\cos\Theta$. 对右端的$\cos(\Theta+\xi)$展开, 平均$\cos\gamma t$为零, $\cos^2\gamma t$为$1/2$, 得
\[
\ddot\Theta+\frac g\ell\sin\Theta
-\frac{a^2\gamma^2}{2\ell^2}\sin\Theta\cos\Theta=0,
\quad
U_{\rm eff}=-mg\ell\cos\Theta+\frac{ma^2\gamma^2}{4}\cos^2\Theta.
\]
取$A=a^2\gamma^2/(2\ell^2)$, 平衡条件为$\sin\Theta(g/\ell-A\cos\Theta)=0$. 向下$\Theta=0$在$A<g/\ell$稳定, 临界$A=g/\ell$时二阶项为零, 有效势增量为$mg\ell\Theta^4/8+O(\Theta^6)$, 在平均模型内仍是非线性稳定的退化最低点; 向上$\Theta=\pi$的有效势曲率始终负, 不稳定. 当$A>g/\ell$时出现两个倾斜稳定平衡$\cos\Theta=g/(\ell A)=2g\ell/(a^2\gamma^2)$, 曲率除以$m\ell^2$为$A\sin^2\Theta>0$. 随强驱动它们趋近水平. 竖直驱动摆具有另一有效势及倒立稳定机制, 不应与本题混淆.
\end{solution}

\begin{exercise}
原题30（选做）. 将$x=x_0\cos\omega_0t+(v_0/\omega_0)\sin\omega_0t=x_1(t,x_0)+x_2(t,v_0)$. 求$a,b,c,d$, 使$x_1(t+\Delta t)=ax_1(t)+bx_2(t)$, $x_2(t+\Delta t)=cx_1(t)+dx_2(t)$.
\end{exercise}
\begin{solution}
官方解答中文改写, 补充退化情况. 设$\delta=\omega_0\Delta t$, 在$x_0v_0\ne0$时, 三角加法给$a=d=\cos\delta$, $b=-x_0\omega_0\sin\delta/v_0$, $c=v_0\sin\delta/(\omega_0x_0)$. 这矩阵依赖选定的初值, 不是一般状态$(x,\dot x)$的传播矩阵.

若$x_0=0$而$v_0\ne0$, $x_1$恒为零, 却无法用仅一个正弦函数表示有余弦项的$x_2(t+\Delta t)$, 除非$\sin\delta=0$. $v_0=0$情形同理. 零解时系数不唯一. 若目标是初值无关的状态传播, 应改用
\[
\binom{x(t+\Delta t)}{\dot x(t+\Delta t)}
=\begin{pmatrix}\cos\delta&\sin\delta/\omega_0\\-\omega_0\sin\delta&\cos\delta\end{pmatrix}
\binom{x(t)}{\dot x(t)}.
\]
这段是解释原题公式的限制, 不将原题悄悄改写成另一题.
\end{solution}

\begin{exercise}
原题31（选做）. 用原题30的$a,b,c,d$组成$M$, 使$\vect x(t+\Delta t)=M\vect x(t)$, 求特征值. 再令$\Delta t=4\pi/\omega_0$求特征向量.
\end{exercise}
\begin{solution}
官方解答中文改写并补全复特征值. 在非退化初值条件下$\det(M-\lambda I)=(\cos\delta-\lambda)^2+\sin^2\delta$, 所以$\lambda_\pm=e^{\pm i\delta}$, 一般为复数. 当$\sin\delta=0$时退化为同一实根$\cos\delta$. 题指定$\delta=4\pi$, 此时$M=I$, 特征值全为1, 每个非零向量都是特征向量; 可取标准基$(1,0)^{\mathsf T},(0,1)^{\mathsf T}$. 不把官方解对实根情形的讨论误读为一般时刻没有特征值.
\end{solution}
